\chapter{Galactorrhea}

\section*{Definition}
Galactorrhea is the spontaneous, persistent secretion of milk-like fluid from the nipple unrelated to pregnancy or breastfeeding. It is usually bilateral and multi-ductal but can be unilateral in some cases.

\section*{What Happens in Your Body}
Galactorrhea results from hyperprolactinemia, which disrupts the normal brain's temperature control center--pituitary inhibition of prolactin (tonic inhibition by dopamine from the tuberoinfundibular pathway). Excess prolactin stimulates lactotroph hyperplasia in the breast and milk production. However, galactorrhea can also occur with normal prolactin levels due to increased breast sensitivity to prolactin.

\section*{Causes}
\begin{itemize}
  \item \textbf{Physiologic:} Pregnancy, breastfeeding, nipple stimulation (mechanical, sexual), newborns (witches' milk), stress, sleep
  \item \textbf{Medication-induced (most common):} Antipsychotics (especially risperidone, haloperidol), metoclopramide, domperidone, SSRIs, verapamil, methyldopa, H2 blockers (cimetidine), opiates, hormonal contraceptives
  \item \textbf{Pituitary/brain's temperature control center:} Prolactinoma (micro- or macroadenoma), pituitary stalk compression, empty sella syndrome, craniopharyngioma
  \item \textbf{Endocrine:} Hypothyroidism (loss of dopamine inhibition from TRH stimulation), polycystic ovary syndrome, chronic renal failure, Cushing disease, acromegaly
  \item \textbf{Chest wall lesions:} Herpes zoster, surgery, trauma, burns (stimulation of thoracic nerves)
\end{itemize}

\section*{When to See a Doctor}
See your doctor if:
\begin{itemize}
  \item Discharge is new, spontaneous, persistent, or associated with headache, visual disturbances, or amenorrhea
  \item Discharge is bloody, unilateral, or from a single duct (may indicate breast pathology)
  \item There are signs of hypogonadism: infertility, loss of libido, erectile dysfunction, oligomenorrhea
\end{itemize}

\section*{Related Symptoms}
\begin{itemize}
  \item Amenorrhea, oligomenorrhea, infertility, loss of libido, erectile dysfunction, gynecomastia, visual field defects, headache, hypothyroidism
\end{itemize}
\textbf{Important:} Galactorrhea is more common in women; in men, hyperprolactinemia more often presents with hypogonadism or mass effects than overt discharge.

\section*{Self-Care / Home Management}
\begin{itemize}
  \item Discontinue causative medications after consulting the prescribing clinician
  \item Minimize breast stimulation (tight clothing, frequent breast exams, sexual stimulation)
  \item For mild cases with normal prolactin and no mass: monitoring may be sufficient
  \item Use breast pads to absorb discharge
\end{itemize}

\section*{How Your Doctor Will Evaluate You}
\begin{itemize}
  \item History: timing, laterality, relation to menses, medications, nipple stimulation
  \item Physical exam: breast exam (mass, discharge character), visual field testing, thyroid palpation
  \item Laboratory: serum prolactin (ensure no recent breast exam or stress), TSH, free T4, renal function, pregnancy test
  \item If hyperprolactinemia confirmed: pituitary MRI with contrast to rule out prolactinoma
  \item Breast imaging (mammogram, ultrasound) if discharge is bloody, unilateral, or single-duct
\end{itemize}

\section*{Treatment Options}
\begin{itemize}
  \item Discontinue offending medications if safe
  \item Treat hypothyroidism with levothyroxine (resolves galactorrhea in most cases)
  \item Dopamine agonists for prolactinoma: cabergoline (preferred) or bromocriptine; normalize prolactin in ~85\% and shrink adenomas
  \item Testosterone replacement for men with hypogonadism secondary to hyperprolactinemia
  \item Surgical resection for large prolactinomas resistant to medical therapy
\end{itemize}

\clearpage
